\paragraph{Scope of validity.} The empirical claims of this paper apply
to cooperative-navigation tasks in the MPE family at agent counts
$N \in \{3, 6, 12\}$, trained with on-policy MAPPO under the
hyperparameters listed in the Experiments section. SMAC and SMACv2
\citep{samvelyan2019smac, ellis2022smacv2} are not evaluated. Neither are
the sparse-reward, partially observable task families --- level-based
foraging and multi-robot warehouse \citep{christianos2020seac} --- on
which communication is load-bearing, nor the vectorised JAX
reimplementations \citep{rutherford2023jaxmarl} that would make a
ten-seed sweep at larger $N$ tractable. Larger $N$ values where the
asymptotic $\mathcal{O}(N^2) \to \mathcal{O}(Nk)$ gap should be most
visible are not exercised in the current run set.

\paragraph{Gate-variance diagnosis is correlational.} The
ninefold inflation of between-seed standard deviation that accompanies
the adaptive-$k$ gate is consistent with gate-induced stochasticity
dominating the policy-gradient signal, but the study does not perform a
controlled intervention isolating the gate as the variance source.
Mechanisms that would establish causality --- freezing the gate to a
fixed-$k$ schedule after pre-training, swapping Gumbel-softmax for a
straight-through Bernoulli, or directly logging the policy-gradient norm
contributed by the gate parameters --- are deferred.

\paragraph{Wall-clock not measured.} The 2080 Ti rollout regime is
CPU-bound, so the FLOPs savings predicted by the
$k_{\mathrm{eff}}/(N{-}1)$ ratio are not realised in measured wall time
at this scale. A GPU-bound rollout regime, which is the relevant regime
for distributed agentic-RL infrastructure, is the natural setting in
which to validate the wall-clock implication of the FLOPs analysis.
